% SAGE Open - 中文草稿（内容审阅用）
% 提交前需替换为英文版本
% 编译说明：本文件含中文，需在文档类添加 ctex 或改用 xelatex 编译，
% 但 SAGE 投稿时须为纯英文。此版本仅供内容审阅。

\documentclass[Afour,sageapa,times]{sagej}

\usepackage{amsmath,amssymb}
\usepackage{booktabs}
\usepackage{float}
\usepackage{enumitem}
\usepackage{array}
\usepackage{multirow}
\usepackage[colorlinks,bookmarksopen,bookmarksnumbered,citecolor=blue,urlcolor=blue]{hyperref}

% === 以下中文内容仅供审阅，正式投稿需翻译为英文 ===

\begin{document}

\runninghead{Xing et al.}

\title{Research on Rapid English Sentence Learning Method Based on\\ Large Language Models and AIGC for Chinese Students}

\author{Xing Shujun\affilnum{1}, Liu Peng\affilnum{1},\\
An Zhengguang\affilnum{1}, Wang Suomin\affilnum{1},\\
and Wei Qinjiang\affilnum{1}}

\affiliation{\affilnum{1}Heze Vocational College, Shandong, China}

\corrauth{Xing Shujun, Heze Vocational College, Heze, Shandong, China.}

\email{xsj@bupt.edu.cn}

\begin{abstract}
【中文摘要】英语句子构造能力是语言产出的核心基础，然而中国学生长期面临"能认词、难造句"的困境。本文提出一种基于大语言模型（LLM）与AIGC的英语句子快速学习法，核心包括：（1）利用AIGC生成15秒电影级短视频，为每个目标单词和例句创建沉浸式情境记忆锚点；（2）构建"理解单词→理解短语→理解句型结构→搞清语法→AI辅助仿写"的五步句子学习法，实现从词汇认知到句子产出的完整转化；（3）通过LLM对学习者仿写句子进行语法正确性、逻辑合理性和表达地道性的三维即时评价与引导性纠错；（4）将艾宾浩斯遗忘曲线嵌入句子复习机制，为正确造句提供自动化间隔重复复习。为验证该方法的有效性，本研究在中国宁夏吴忠地区开展为期12周的三组对照实验（$n=141$）。结果显示，实验组在句子构造准确率各维度上均显著优于传统教学组和传统学习机组（$p<0.001$），表达地道性的效应量最大（Cohen's $d=1.28$），延迟后测词汇保持率达80.5\%。本研究为AI技术赋能英语句子教学提供了可复制的实践方案。
\end{abstract}

\keywords{large language model (LLM), AI-generated content (AIGC), English sentence learning, real-time evaluation, Ebbinghaus forgetting curve}

\maketitle


% ============================================================
\section{Introduction}  % 引言

英语句子构造能力是衡量语言综合运用能力的关键指标。然而，中国学生长期面临"能认词、难造句"的困境。Li等\cite{li2025twoyears}对144篇2023--2024年间生成式AI语言学习研究的系统综述表明，现有研究以写作能力为主（占42.4\%），对句子产出训练的专项研究仍然不足。中国95\%以上英语教学仍采用传统课堂模式，在农村地区这一比例更高。传统课堂以词汇背诵和语法讲解为主，班额通常超过50人\cite{wang2011large}，学生缺乏充分的句子产出练习——一个40人的班级，每节课能够获得教师书面批改造句反馈的学生通常不超过5人，反馈延迟可达数天之久。当前中国市场上的英语学习软件普遍仅关注词汇的"音—形—义"识别，缺乏对句子层面输出的训练和评价。

大语言模型（LLM）的成熟为句子教学带来了革命性突破。以GPT系列、DeepSeek、通义千问等为代表的LLM能够对学习者提交的句子进行秒级、多维度的评价与纠错——涵盖语法正确性、逻辑合理性和表达地道性三个维度。Mahapatra\cite{mahapatra2024impact}在一项混合方法干预研究中证实，ChatGPT辅助的写作教学显著提升了ESL学生的写作技能；Hao和Razali\cite{hongxia2025impact}针对中国EFL本科生的实验进一步表明，ChatGPT不仅提高了英语学术写作能力，还显著增强了学习者的参与度和自主性。LLM使"造句→即刻反馈→修正→再反馈"的高频学习循环成为可能，从根本上突破了传统教学中反馈延迟和频次不足的瓶颈。

与此同时，AIGC视频生成技术的突破使得为每个单词和例句定制电影级情境视频成为可能。目前生成一段1秒电影级情节短视频的成本已降至约1元人民币——仅为传统影视制作的千分之一。Chen等\cite{chen2024empowering}对134篇AIGC教育应用文献的系统综述指出，AIGC正在成为教育变革的重要驱动力，尤其在多媒体学习资源生成方面展现出巨大潜力。电影级短视频将目标单词及其例句嵌入一个完整的微型叙事中，激活学习者的情境记忆（Episodic Memory），建立"情境-单词-例句"的三维记忆锚点。需要指出的是，本方法中的AIGC专指基于扩散模型的视频生成技术，用于创建情境视频记忆锚点；而LLM承担独立的自然语言评价与纠错功能。两者属于不同的AI技术分支，在本方法中协同工作但不可相互替代。

此外，Huang\cite{huang2025spaced}从认知心理学视角系统论证了间隔重复与提取练习的有效性，并探讨了AI如何赋能这两种学习策略，为将艾宾浩斯遗忘曲线与现代智能学习系统整合提供了理论框架。

本文的核心贡献在于：（1）首次系统性地提出"理解单词→理解短语→理解句型结构→搞清语法→AI辅助仿写"的五步句子学习法，形成从词汇认知到句子产出的完整链路；（2）将AIGC电影级短视频、LLM即时评价和艾宾浩斯遗忘曲线复习整合为句子学习的闭环系统；（3）通过严格的12周三组对照实验，在真实教学场景中验证了该方法在句子构造准确率方面的显著优势。


% ============================================================
\section{Literature Review}  % 文献综述

\subsection{Sentence Production and Language Acquisition}

Swain的输出假设（Output Hypothesis）指出，语言产出在二语习得中发挥着不可替代的作用——学习者在尝试产出目标语言句子时，会注意到自身语言能力与表达需求之间的差距（Noticing Gap），这种认知冲突驱动了语言系统的发展\cite{swain1995three}。然而，在中国占据主导地位的传统课堂教学中，句子产出练习频率低、反馈延迟大，严重制约了这一习得机制的发挥。提高句子产出的频率和反馈质量是提升二语能力的关键。

\subsection{AIGC Situational Video and Episodic Memory}

多媒体学习认知理论（Cognitive Theory of Multimedia Learning, CTML）为视频辅助语言学习提供了理论基石。Mayer在2024年对该理论进行了全面回顾与展望，指出多媒体学习的核心在于以双通道（视觉+听觉）呈现信息时能够分散认知负荷、促进深层加工\cite{mayer2024past}。Tulving的情境记忆理论进一步指出，与特定时间、地点和情感体验绑定的记忆比纯粹的语义记忆具有更强的提取效率\cite{tulving1983elements}。Paivio的双重编码理论阐释了言语编码与意象编码双通道在记忆提取中的协同优势\cite{paivio1991dual}。

在教育视频研究方面，Navarrete等\cite{navarrete2025closer}对257篇视频学习文献的系统综述表明，视频特征（如时长、叙事结构、互动性）对学习效果有显著调节作用。Monib等\cite{khan2025microlearning}对40项微学习研究的系统综述进一步确证：时长在数秒至数分钟的"精炼内容"以目标驱动、动作导向为特征，在认知（知识保持、迁移）、行为（参与度、完成率）和情感（动机、满意度）三个维度上均显著优于传统长时教学。

在AI生成的视觉内容领域，Behforouz和Al Ghaithi\cite{behforouz2025visual}通过随机对照实验发现，利用AI工具生成动画辅助词汇学习的学生，其词汇增长（$p<0.05$）和学习动机均显著优于对照组。Zhang和Zhang\cite{zhang2024multimedia}的实验研究进一步表明，多媒体增强的词汇输入条件（口语+视觉+副语言线索）对词汇习得效果有显著促进作用。

综合上述研究框架，AIGC电影级短视频的设计原则可归纳为：（1）双通道编码——视觉场景+听觉语音同时激活双重编码路径；（2）微叙事原则——15秒微型叙事具有完整的起承转合，在认知负荷可控的范围内建立情境记忆锚点；（3）情感激活——电影级质感与戏剧张力触发情感记忆，增强编码深度。

\subsection{LLM in Language Education}

大语言模型在语言教育中的应用是当前研究热点。Kasneci等指出，LLM可用于个性化辅导、作业批改和自适应学习路径规划\cite{kasneci2023chatgpt}。Yan和Zhang\cite{yan2024l2}通过混合方法多案例研究发现，L2学习者对ChatGPT提供的自动书面纠错反馈（AWCF）表现出积极的行为参与和认知参与，LLM反馈在语法准确性、词汇适用性和篇章连贯性方面均展现出显著潜力。然而，现有研究多集中在文本生成和阅读理解辅助方面，将LLM作为句子产出的秒级三维评价引擎、并整合AIGC视频和艾宾浩斯复习系统的完整方法研究尚属空白。

\subsection{Ebbinghaus Forgetting Curve and Sentence Review}

德国心理学家Ebbinghaus在1885年提出的遗忘曲线揭示了人类记忆随时间指数衰减的规律\cite{ebbinghaus1885memory}。间隔重复策略已被广泛应用于词汇记忆领域，但将完整句子纳入间隔重复系统的研究尚不多见。本文提出的方法将学习者仿写正确的句子纳入复习队列，在遗忘临界点自动推送复习任务，实现了句子层面的间隔重复巩固。


% ============================================================
\section{Core Methodology}  % 核心方法

本方法的核心教学流程是"AI辅助句子五步学习法"，每一步都有AI参与学生的学习与评价，每一步都以前一步为基础，层层递进，确保从单词认知到句子产出的完整转化。在本方法中，单词真正掌握的标志是能够造出正确的句子，句子真正掌握的标志是能够仿写出正确的句子。传统教学环境下，学生难以独立完成这一闭环，而AI的介入使得每一步均可获得即时辅助与反馈，从而打通从词汇识别到独立造句的完整通路。

\subsection{Step 1: Understanding Word Meaning and Usage}
系统可通过AI交互使用户学会单词的音标、中文释义、词性和例句，AI提供详细解释，不懂的就降低难度。这一阶段的目标是在AI辅助下，从"知道有这个单词"上升到"理解其使用场景"。学习者观看AIGC为该单词生成的电影级短视频，在情境中初次接触目标单词。

系统为每个目标单词和例句利用AI生成专属的电影级短视频。EnglishUp学习平台利用AIGC技术，根据提示词自动生成15秒以内的电影级情节短视频，遵循"微叙事"原则：0--5秒建立场景氛围，5--10秒聚焦角色情感，10--15秒呈现目标例句。例如，对于单词``reluctant''及其例句``I'm reluctant to leave the party.''，提示词设计为：

\begin{quote}\small
\textit{16:9 cinematic 15s English vocabulary teaching video, warm soft lighting, shallow depth of field, film-grade texture. Scene: A little boy plays blocks happily at a colorful kids' party. His mother walks over hurriedly with a schoolbag and urges: ``We have class this afternoon. Pack up quickly.'' The boy holds his toy tightly, looks reluctant, and says softly: ``I'm reluctant to leave the party.'' Requirements: Chinese-English bilingual subtitles, key word ``reluctant'' highlighted.}
\end{quote}

生成的15秒视频为学生创造了一个多重感官参与的记忆编码环境，初次学习时用于情境导入，复习时用于记忆唤醒。

\subsection{Step 2: Understanding Collocations through AI}
单词很少孤立使用，掌握其常见搭配是从"认识单词"到"使用单词"的关键。AI系统为每个目标单词自动提取3--5个高频短语或搭配，每个配以AIGC生成的简短场景视频（5--8秒）和AI的用法讲解。例如，对于单词``impression''，AI展示``make a good impression''、``first impression''、``leave a deep impression''，并配以相应场景视频。

\subsection{Step 3: Understanding Sentence Structures through AI}
AI自动分析目标单词所在的原句，为学习者拆解句型结构和语法框架。本方法建立了分级句型库：小学阶段覆盖200个核心句型（如\textit{There be...}、\textit{I want to...}），初中阶段覆盖300个核心句型（如宾语从句、被动语态等）。AI通过句型拆解和类比，帮助学习者理解句子结构规律。

\subsection{Step 4: Clarifying Grammar through AI}
AI系统针对当前句子的语法要点、易错点和知识盲区进行精准诊断——如时态混淆、介词搭配错误、词性误用等——生成针对性讲解和预警，为自主句子产出扫清障碍。

\subsection{Step 5: AI-Assisted Imitation Writing}
仿写是检验句子学习效果的最终环节。学习者使用目标单词和句型独立仿写3--5个原创句子（而非简单重复例句）。系统将每个句子提交给LLM进行三维评价：（a）语法正确性，检查主谓一致、时态、冠词等要素；（b）逻辑合理性，判断语义自洽性；（c）表达地道性，评估是否符合母语者习惯。如存在错误，AI提供引导性提示而非直接答案，帮助学习者自主修正，过程重复至句子达标。

LLM即时评价机制实现了四项核心功能：秒级反馈（1--3秒内返回结果）、多维度诊断（如"第三人称单数未加-s"）、脚手架式纠错（提示→引导→示范）和正向激励。

\subsection{Ebbinghaus Review System}

所有通过第五步AI评价达标的句子，自动进入艾宾浩斯遗忘曲线复习系统，复习时间节点为学习后的第1天、第2天、第4天、第7天、第15天和第30天。每次复习时系统自动推送包含原AIGC视频的复习任务，要求学习者再次造句、AI再次评价。若复习失败（造句出错），降级至更短间隔；连续三次通过则适当延长。未掌握的句子可标记为"收藏"，始终保持在复习队列中。


% ============================================================
\section{Experimental Design}  % 实验设计

\subsection{Experimental Platform}
本实验的实验平台为自研的\textbf{EnglishUp}智能英语学习系统，是一套经过多年研发和迭代的成熟的AI英语学习软件。EnglishUp集成了AIGC电影级短视频生成、AI辅助句子五步学习法、LLM即时评价与纠错、艾宾浩斯遗忘曲线复习等核心功能，前端通过平板设备为学生提供统一的学习界面，支持离线缓存与在线同步。实验组学生统一配发搭载EnglishUp系统的平板设备。

\subsection{Participants}
本实验在宁夏吴忠地区的一所初级中学开展。从八年级6个平行班中随机抽取3个班作为实验对象。所有参与学生均在其英语晚辅课期间使用对应方法进行实验，三组晚辅时长相同，且保证各组在晚辅时段和常规上课时段之外不进行额外的英语词汇或句子训练。

\begin{table}[H]
\centering
\small
\caption{Basic Information of Participants}
\begin{tabular}{lcccc}
\toprule
\textbf{Group} & \textbf{N} & \textbf{Male (\%)} & \textbf{Pretest ($M \pm SD$)} & \textbf{Method} \\
\midrule
A (Control 1) & 48 & 52.1 & 68.3 $\pm$ 12.5 & Traditional classroom \\
B (Control 2) & 46 & 50.0 & 67.9 $\pm$ 13.1 & Ebbinghaus vocab device\tnote{a} \\
C (Experimental) & 47 & 51.1 & 68.1 $\pm$ 12.8 & AI five-step sentence method \\
\bottomrule
\end{tabular}
\begin{tablenotes}
\item[a] Representative of traditional vocabulary learning devices widely available in the Chinese consumer market, which use Ebbinghaus-based spaced repetition for word and sentence memorization only, without AI sentence evaluation or AIGC video features.
\end{tablenotes}
\end{table}

三组前测无显著差异（$F(2,138)=0.02$, $p=0.98$），基线一致。

\subsection{Experimental Procedure and Evaluation Metrics}
实验周期为12周（准备期第1周、实验期第2--11周、评估期第12周），所有组别学习700个中考考纲核心词汇。三组均利用每天英语晚辅课时段（时长相同）进行学习：实验组通过AI辅助句子五步学习法学习3--5个句子和10个以上新单词；B组使用中国市场大量存在的传统学习机以艾宾浩斯算法背3--5个句子和10个以上单词（该类设备提供词汇和句子机械记忆功能，无AI句子评价与AIGC视频功能）；A组按传统课堂进度讲授（代表中国95\%以上英语课堂的典型教学模式，农村地区占比更高）。晚辅时段和常规上课时段之外，三组均不进行额外的英语词汇或句子训练。核心评估指标为句子构造准确率（SCA），由两位英语教师按三维标准盲审评分（5分制，评分者间信度 Cohen's $\kappa = 0.86$），辅以词汇记忆保持率（VRR）和学习效率（LE）。其中，VRR通过听写测试评估：从700个目标词汇中随机抽取100词，要求学生在规定时间内写出对应的英文单词及中文释义，正确率记为VRR得分。即时后测于第12周实验结束时进行，延迟后测于第14周（即实验结束后2周）进行，两次测试采用等值但不同题目的平行试卷。学习效率（LE）定义为单位学习时间内有效掌握的词汇数量。


% ============================================================
\section{Experimental Results and Analysis}  % 实验结果与分析

\subsection{Sentence Construction Accuracy}
句子构造准确率是本研究的核心指标。三组的SCA得分如表\ref{tab:sca}所示。

\begin{table}[H]
\centering
\small
\caption{Comparison of Sentence Construction Accuracy Across Three Groups\label{tab:sca}}
\begin{tabular}{lccc}
\toprule
\textbf{Group} & \textbf{Grammar} & \textbf{Logic} & \textbf{Nativeness} \\
\midrule
A (Traditional) & 3.18 $\pm$ 0.87 & 2.95 $\pm$ 0.93 & 2.42 $\pm$ 0.84 \\
B (Device) & 3.62 $\pm$ 0.76 & 3.28 $\pm$ 0.85 & 2.75 $\pm$ 0.80 \\
C (Experimental) & \textbf{4.28 $\pm$ 0.52} & \textbf{4.02 $\pm$ 0.60} & \textbf{3.72 $\pm$ 0.65} \\
\bottomrule
\end{tabular}
\end{table}

注：评分采用5分制（1=完全错误，5=完全正确）。Grammar = 语法正确性，Logic = 逻辑合理性，Nativeness = 表达地道性。

C组在所有三个维度上均显著领先（$p<0.001$）。语法正确性：$F(2,138)=28.7$, $\eta^2=0.29$, C组 vs B组 Cohen's $d=0.96$；逻辑合理性：$F=19.4$, $\eta^2=0.22$, $d=0.93$；表达地道性：$F=35.6$, $\eta^2=0.34$, $d=1.28$。表达地道性的效应量最大（$d=1.28$），表明AIGC视频的母语者场景输入和LLM对地道性的逐句打磨，显著提升了学习者句子产出的自然度。

\subsection{Vocabulary Retention Rate}

\begin{table}[H]
\centering
\small
\caption{Comparison of Vocabulary Retention Rates Across Three Groups}
\begin{tabular}{lccc}
\toprule
\textbf{Group} & \textbf{Immediate (\%)} & \textbf{Delayed--2 wk (\%)} & \textbf{Decay (\%)} \\
\midrule
A & 64.8 $\pm$ 9.2 & 44.1 $\pm$ 10.8 & -20.7 \\
B & 80.1 $\pm$ 7.5 & 64.3 $\pm$ 8.8 & -15.8 \\
C & \textbf{89.2 $\pm$ 6.1} & \textbf{80.5 $\pm$ 7.3} & \textbf{-8.7} \\
\bottomrule
\end{tabular}
\end{table}

方差分析显示$F(2,138)=52.6$, $p<0.001$, $\eta^2=0.43$。C组衰减仅8.7个百分点，远低于A组20.7和B组15.8个百分点。句子层面的深度加工对词汇记忆有显著的"以句带词"促进作用。

\subsection{Learning Efficiency and Statistical Summary}
C组总学习时长50小时，词汇即时掌握率89.2\%（按100词抽样外推，约合624词，LE=12.48词/小时）；B组LE=11.2（纯刷词模式，无句子产出训练）；A组LE=9.10。C组在句子产出能力和词汇学习效率上均领先于B组。所有指标三组差异均达$p<0.001$水平，事后检验确认在句子产出能力维度上C组显著优于B组和A组，形成明确的阶梯效应。


% ============================================================
\section{Discussion}  % 讨论

\subsection{Mechanisms Underlying Sentence Production Improvement}
本研究结果表明，AI辅助句子五步学习法的优势可从四个机制解释：

\textbf{（1）多模态情境编码驱动深度理解。} AIGC视频为每个单词和例句创造了独特的视觉—听觉—情感记忆标签，使得抽象的语法规则和词汇在具体的剧情场景中获得语义锚点，为句子产出提供丰富的提取线索，这与双重编码理论\cite{paivio1991dual}和情境记忆理论\cite{tulving1983elements}的核心预测一致。

\textbf{（2）输出驱动深度加工。} 仿写环节迫使学习者从单词"再认"上升到句子"产出"，将被动句法知识转化为主动产出能力。Swain的输出假设指出，每轮"造句→评价→修正"都是一次"假设—验证"过程\cite{swain1995three}。

\textbf{（3）即时反馈加速学习循环。} LLM评价将传统"造句→等待批改→收到反馈→修正"的漫长循环压缩至秒级。实验期间，实验组学生人均完成超过600次AI评价，这是传统教学无法实现的反馈频次。

\textbf{（4）句子层面的间隔复习。} 将完整句子纳入艾宾浩斯系统，使每轮复习强化整个句法结构而不仅是孤立词汇。

\subsection{Limitations}
本研究的局限性包括：（1）12周的实验周期有限，方法的长期效果仍需追踪；（2）主要聚焦书面句子产出，口语造句能力尚未评估。


% ============================================================
\section{Conclusion and Future Work}  % 结论与展望

本文基于自研EnglishUp平台，提出并验证了AI辅助句子五步学习法。实验组在语法正确性（4.28）、逻辑合理性（4.02）和表达地道性（3.72）上均显著领先（$p<0.001$），地道性效应量$d=1.28$，延迟记忆保持率达80.5\%。研究表明，AI辅助句子五步学习法+艾宾浩斯句子复习的整合方法能够有效提升学习者的英语句子产出能力。未来将拓展至口语句子训练，并在更广泛的地区开展验证。


% ============================================================
% Data Availability
\section*{Data Availability Statement}
本研究的部分示例实验数据已在GitHub开源公开：\url{https://github.com/xsj1984/EnglishUPbyAIGC}。本研究自研的实验平台网址：\url{https://engup.cn}。

% ============================================================
% Declarations
\begin{acks}
本研究项目经费来源于作者邢树军个人投入（120万元人民币），旨在为解决中国农村学生英文教育问题提供AI赋能的公益实践。实验数据独立于EnglishUp平台商业运营，实验过程由第三方教师监督。
\end{acks}

\begin{funding}
本项目由作者邢树军个人资助（120万元人民币），致力于通过AI技术改善中国农村地区英语教育水平。
\end{funding}

\begin{dci}
作者声明无利益冲突。
\end{dci}

\section*{Author Contributions}
邢树军（Xing Shujun）：研究构思、方法论设计、EnglishUp平台开发、实验实施、初稿撰写、项目资助；刘鹏（Liu Peng）：实验数据采集与整理、统计分析；安政光（An Zhengguang）：实验实施、教学资源准备；王所民（Wang Suomin）：实验实施、数据采集；魏钦江（Wei Qinjiang）：平台技术支持、数据分析。所有作者均已阅读并同意最终版本。

\section*{Ethical Approval}
本研究为常规课堂教学实验，不涉及任何伦理风险。实验前已获得所有参与学生及其监护人的知情同意，实验数据严格匿名化处理。


% ============================================================
% References (APA 7th, to be converted from current Chinese format)
\begin{thebibliography}{99}

\bibitem{li2025twoyears}
Li, B., Tan, Y. L., Wang, C., \& Lowell, V. (2025). Two years of innovation: A systematic review of empirical generative AI research in language learning and teaching. \textit{Computers and Education: Artificial Intelligence}, 9, 100445.

\bibitem{wang2011large}
Wang, Q. (2011). \textit{Teaching in large classes in China: Challenges and opportunities}. University of Warwick.

\bibitem{mahapatra2024impact}
Mahapatra, S. (2024). Impact of ChatGPT on ESL students' academic writing skills: A mixed methods intervention study. \textit{Smart Learning Environments}, 11, 9.

\bibitem{hongxia2025impact}
Hao, H., \& Razali, A. B. (2025). Impact of ChatGPT on English academic writing ability and engagement of Chinese EFL undergraduates. \textit{International Journal of Instruction}, 18(2), 323--346.

\bibitem{chen2024empowering}
Chen, X., Hu, Z., \& Wang, C. (2024). Empowering education development through AIGC: A systematic literature review. \textit{Education and Information Technologies}, 29, 17485--17537.

\bibitem{huang2025spaced}
Huang, M. (2025). Spaced repetition and retrieval practice: Efficient learning mechanisms from a cognitive psychology perspective and their empowerment by AI. \textit{International Journal of Asian Social Science Research}.

\bibitem{swain1995three}
Swain, M. (1995). Three functions of output in second language learning. In G. Cook \& B. Seidlhofer (Eds.), \textit{Principle and Practice in Applied Linguistics} (pp. 125--144). Oxford University Press.

\bibitem{mayer2024past}
Mayer, R. E. (2024). The past, present, and future of the cognitive theory of multimedia learning. \textit{Educational Psychology Review}, 36, Article 8.

\bibitem{tulving1983elements}
Tulving, E. (1983). \textit{Elements of Episodic Memory}. Oxford University Press.

\bibitem{paivio1991dual}
Paivio, A. (1991). Dual coding theory: Retrospect and current status. \textit{Canadian Journal of Psychology}, 45(3), 255--287.

\bibitem{navarrete2025closer}
Navarrete, E., Nehring, A., Schanze, S., Ewerth, R., \& Hoppe, A. (2025). A closer look into recent video-based learning research: A comprehensive review of video characteristics, tools, technologies, and learning effectiveness. \textit{International Journal of Artificial Intelligence in Education}, 35, 1631--1694.

\bibitem{khan2025microlearning}
Monib, W. K., Qazi, A., \& Apong, R. A. (2025). Microlearning beyond boundaries: A systematic review and a novel framework for improving learning outcomes. \textit{Heliyon}, 11(2), e41413.

\bibitem{behforouz2025visual}
Behforouz, B., \& Al Ghaithi, A. (2025). Visual language learning: Exploring the effects of artificial intelligence-generated animations on vocabulary and motivational growth. \textit{Technology in Language Teaching \& Learning}, 7(4), Article 103077.

\bibitem{zhang2024multimedia}
Zhang, P., \& Zhang, S. (2024). Multimedia enhanced vocabulary learning: The role of input condition and learner-related factors. \textit{System}, 122, 103275.

\bibitem{kasneci2023chatgpt}
Kasneci, E., et al. (2023). ChatGPT for good? \textit{Learning and Individual Differences}, 103, 102274.

\bibitem{yan2024l2}
Yan, D., \& Zhang, S. (2024). L2 writer engagement with automated written corrective feedback provided by ChatGPT: A mixed-method multiple case study. \textit{Humanities and Social Sciences Communications}, 11, Article 1086.

\bibitem{ebbinghaus1885memory}
Ebbinghaus, H. (1885). \textit{Memory: A Contribution to Experimental Psychology}. Teachers College, Columbia University.

\end{thebibliography}

\end{document}
